\documentclass[aps,prl,twocolumn,superscriptaddress]{revtex4-2}
\usepackage{amsmath,amssymb,graphicx,hyperref}
\begin{document}

\title{Replication Report: \emph{Magnetic control of spin-orbit fields:
a first-principles study of Fe/GaAs junctions}}
\author{Replication (subagent \texttt{texture-gmitra2013})}
\date{\today}

\begin{abstract}
We replicate the symmetry-based spin-orbit-field (SOF) framework of Gmitra,
Matos-Abiague, Draxl and Fabian [PRL \textbf{111}, 036603 (2013);
arXiv:1303.2510]. The paper's DFT (relativistic FLAPW/FLEUR on Fe/GaAs slabs) is
out of scope for in-process replication; instead we reimplement the paper's
analytic core---the $C_{2v}$ SOC Hamiltonian, the Bychkov--Rashba/Dresselhaus
decomposition, the symmetry extraction formulas (Eqs.~4--9), and the
magnetization-angle parametrization of Table~I---and verify the internal
consistency and the downstream physical claims. All six machine-checkable claims
pass. In particular we reproduce (i) the exact self-consistency of the extraction
method, (ii) the sign change of $\alpha_1(\theta)$, (iii) the sign flip of
$\alpha_1\beta_1$ between $[1\bar10]$ and $[110]$ (band $n{=}2$ does not flip),
(iv) the linear-in-$k$ regime and anisotropic ``butterfly'' textures, and
(v) the pure $\cos(2\theta)$ $C_{2v}$ angular form with band~1 far more
magnetization-sensitive than band~2.
\end{abstract}
\maketitle

\section{Central claims}
The paper's headline results are:
\begin{enumerate}
\item SOFs at the $C_{2v}$ Fe/GaAs(001) interface depend not only on the
interface electric field but \emph{surprisingly strongly on the Fe
magnetization orientation}---enabling magnetic control.
\item The k-space SOF patterns are highly anisotropic and change qualitatively
with band/energy (``butterflies''); the anisotropy axes can be \emph{flipped}
by rotating the magnetization.
\item A symmetry-based method extracts $\mathbf{w}(\mathbf{k})$ (magnitude and
direction) for a generic Bloch state directly from ab-initio bands, valid also
at large $k$ (not just near $\Gamma$).
\end{enumerate}

\section{Model}
The most general $C_{2v}$-allowed in-plane SOC Hamiltonian (Eq.~1) is
\begin{equation}
H_{\mathrm{so}}=\mu_n(k_x,k_y,\theta)\,k_x\sigma_y
+\eta_n(k_x,k_y,\theta)\,k_y\sigma_x,
\end{equation}
with $x\parallel[1\bar10]$, $y\parallel[110]$, $\theta$ the magnetization angle
from $[1\bar10]$, and (Eq.~2) $\mu_n=\mu_n^{(0)}+\mu_n^{(1)}k_x^2+\dots$,
$\eta_n=\eta_n^{(0)}+\eta_n^{(1)}k_x^2+\dots$. Writing $H_{\mathrm{so}}=
\mathbf{w}_n(\mathbf{k})\cdot\boldsymbol{\sigma}$ (Eq.~3),
$\mathbf{w}_n=(\eta_n k_y,\ \mu_n k_x,\ 0)$. Near $\Gamma$,
$\alpha_n=(\mu_n^{(0)}-\eta_n^{(0)})/2$ (Bychkov--Rashba) and
$\beta_n=(\mu_n^{(0)}+\eta_n^{(0)})/2$ (Dresselhaus). The angular dependence
(Table~I, Eqs.~10--11) is
\begin{align}
\alpha_n(\theta)&\simeq A_n^{(+)}+B_n^{(+)}\cos(2\theta),\\
\beta_n(\theta)&\simeq A_n^{(-)}+B_n^{(-)}\cos(2\theta),
\end{align}
with the paper's Table~I coefficients (meV\,\AA) used verbatim.

The symmetry extraction (Eqs.~4--7) reads
$w_{nx}=\sigma\,[\Delta E_n^{\mathrm{so}}+\Gamma_n^{\mathrm{so}}]/(2\cos\theta)$,
$w_{ny}=\sigma\,[\Delta E_n^{\mathrm{so}}-\Gamma_n^{\mathrm{so}}]/(2\sin\theta)$,
with $\Delta E_n^{\mathrm{so}}=[E_n(\mathbf{k})-E_n(-\mathbf{k})]/2$ and
$\Gamma_n^{\mathrm{so}}=[E_n(-k_x,k_y)-E_n(k_x,-k_y)]/2$; the near-$\Gamma$
parameters follow from band velocities $a_n,b_n$ via Eqs.~(8--9).

\section{Method (this work)}
We implement the above in \texttt{code/sof\_model.py} and verify claims in
\texttt{code/run\_analysis.py}. We build a toy band
$E_n(\mathbf{k},\theta)=c(k_x^2+k_y^2)+\mathbf{w}_n(\mathbf{k})\cdot\hat{\mathbf m}$
(first-order in-plane perturbation theory, $\hat{\mathbf m}=(\cos\theta,
\sin\theta)$) and feed it back through Eqs.~(4--9) to confirm the extraction is
exact. We take Table~I as the ab-initio ground truth and check the physical
consequences the paper draws from it.

\section{Results}
All results in \texttt{results/metrics.json}; figures in \texttt{figs/}.
\begin{itemize}
\item \textbf{C1 (extraction round-trip).} Eqs.~(4--9) recover the input
$(\alpha,\beta)$ to $<10^{-6}$ meV\,\AA{} and $w_{x,y}$ to $<10^{-9}$ meV across
bands $n{=}1$--$4$ and nine test angles. The method is exact by construction.
\item \textbf{C2 ($\alpha_1$ sign change).} $\alpha_1(\theta)\in[-6.68,+5.84]$
meV\,\AA{} crosses zero: magnetic control confirmed.
\item \textbf{C3 (product sign flip).} $\alpha_1\beta_1$: $+46.8$ at
$[1\bar10]$, $-91.3$ at $[110]$ (flips). $\alpha_2\beta_2$: $+2419$ vs $+2502$
(no flip). Matches the paper exactly---band~1 axes flip, band~2 preserved.
\item \textbf{C4 (linearity + anisotropy).} $|w|/k$ is $k$-independent in the
linear regime (rel.\ dev.\ $<10^{-9}$). On the $k=\pi/8d$ contour the max/min
of $|w|/k$ is $42.6$ ($M\parallel[1\bar10]$) vs $2.19$ ($M\parallel[110]$):
strongly anisotropic butterflies whose shape changes with $M$.
\item \textbf{C5 ($C_{2v}$ angular form).} Angular modulation $|B/A|$ for
$\alpha$: band~1 $=14.9$ vs band~2 $=0.043$---band~1 is $\sim350\times$ more
magnetization-sensitive, matching the paper's statement.
\end{itemize}

\section{Figures}
\texttt{fig3\_soc\_vs\_theta.png} reproduces the E$=0$ structure of Fig.~3
($\alpha_n,\beta_n$ vs $\theta$). \texttt{fig2\_sof\_butterflies.png} and
\texttt{fig2\_polar\_wk.png} reproduce the Fig.~2 SOF vector fields and polar
$|w|/k$ on the $\pi/100d,\pi/8d,\pi/5d$ contours for band~1 at
$M\parallel[1\bar10]$ and $[110]$.

\section{Verdict}
\textbf{Partial--to--strong replication of the analytic core.} The DFT-derived
numbers themselves are not reproduced (out of scope), but every symmetry claim
and every downstream physical consequence the paper draws from its extracted
Table~I is reproduced exactly. \textbf{Coverage: 6/10} (analytic core + all
qualitative claims covered; ab-initio band structure, electric-field sweep, and
out-of-plane $\langle s\rangle/\Delta_{xc}$ maps not reproduced).
\textbf{Agreement: 9/10} (internal consistency and all sign/anisotropy/ordering
claims match; capped below 10 because the DFT inputs are taken as given rather
than independently computed).

\end{document}
